
\section{\hard\ \trSerre{}对\trAmpleness{}与\trAffineness{}的刻画}

\mbox{定理~16.2.2} 曾从射影几何、\trGlobalGeneration{}与\trZariskiTopology{}等角度给出\trAmpleness{}的若干刻画。这里是又一条刻画——这一次是上同调的——在\trNoetherian{}假设之下。由于（有点出人意料地）我们不会用到这个结果，本节标以星号。

\textbf{18.8.1.\ \trhTheorem（\trAmpleness{}的\trSerre{}上同调判别法）。}\ ——\ \emph{设 $A$ 是一个\trNoether{}环，$X$ 是一个\trProper{} $A$-\trScheme，而 $\mathscr{L}$ 是 $X$ 上一个\trInvertibleSheaf。则以下诸条件等价。}

（a--c）\emph{\trInvertibleSheaf{} $\mathscr{L}$ 在 $X$ 上是\trAmple{}的（在 $A$ 上\trAmple）。}

（e）\emph{对 $X$ 上一切\trCoherentSheaf{} $\mathscr{F}$，存在某个 $n_0$，使得当 $n\ge n_0$ 时，对一切 $i>0$ 都有 $\mathrm{H}^i(X,\mathscr{F}\otimes\mathscr{L}^{\otimes n})=0$。}

标签 $(a)$--$(c)$ 的用意是呼应\mbox{定理~16.2.2} 的陈述。我们避开标签 $(d)$，因为它已被\mbox{定理~16.2.6} 用掉了。（题外话：「\trProperness」这一假设是去不掉的；见\mbox{例~18.7.6}。）在进入证明之前，我们先给出一批应用，以此说明这个结果的动机。（作为热身，你可以给出\mbox{练习~16.2.G} 在\trNoetherian{}情形下的第二种解法：利用 $\pi$ 的\trAffineness{}证明 $\mathrm{H}^i(X,\mathscr{F}\otimes(\pi^*\mathscr{L})^{\otimes m})=\mathrm{H}^i(Y,(\pi_*\mathscr{F})\otimes\mathscr{L}^{\otimes m})$。）

\textbf{18.8.A.（练习）}\quad 设 $X$ 是一个\trProper{} $A$-\trScheme（$A$ \trNoetherian），而 $\mathscr{L}$ 是 $X$ 上一个\trInvertibleSheaf。证明：$\mathscr{L}$ 在 $X$ 上是\trAmple{}的，当且仅当 $\mathscr{L}|_{X^{\mathrm{red}}}$ 在 $X^{\mathrm{red}}$ 上是\trAmple{}的。提示：「仅当」方向用\mbox{练习~16.2.G}。「当」方向：设 $\mathscr{I}$ 为在 $X$ 中切出\trClosedSubscheme{} $X^{\mathrm{red}}$ 的那个\trIdealSheaf。把 $\mathscr{F}$ 按 $\mathscr{I}$ 的幂过滤：
\[
0=\mathscr{I}^r\mathscr{F}\subset\mathscr{I}^{r-1}\mathscr{F}\subset\cdots\subset\mathscr{I}\mathscr{F}\subset\mathscr{F}.
\]
（本质上同一个滤过曾出于类似的理由出现在\mbox{练习~18.4.L} 中。）证明：每个商 $\mathscr{I}^n\mathscr{F}/\mathscr{I}^{n-1}\mathscr{F}$ 在用 $\mathscr{L}$ 的一个足够高的幂扭转之后，都没有更高的\trCohomology。对 $n$ 作\trDescendingInduction，证明滤过的每一段 $\mathscr{I}^n\mathscr{F}$（从而特别是 $\mathscr{F}$）也具有这条性质。

\textbf{18.8.B.（练习）}\quad 设 $X$ 是一个\trProper{} $A$-\trScheme（$A$ \trNoetherian），而 $\mathscr{L}$ 是 $X$ 上一个\trInvertibleSheaf。证明：$\mathscr{L}$ 在 $X$ 上是\trAmple{}的，当且仅当 $\mathscr{L}$ 在每个分量上都\trAmple。提示：仿照上一道练习之解答的概要来做，改取 $\mathscr{I}$ 为其中一个分量的\trIdealSheaf。或许宜先化归到 $X=X^{\mathrm{red}}$ 的情形。

\textbf{18.8.C.（棘手练习）}\quad 设 $C$ 是一条\trProper{}、\emph{\trReduced}的\trCurve（在域 $k$ 上），而 $\mathscr{L}$ 是 $C$ 上一个\trLineBundle，使得 $\nu^*\mathscr{L}$ 是 $\widetilde{C}$ 上的一个\trAmple{}\trLineBundle（正规化一如既往是 $\nu\colon\widetilde{C}\to C$）。证明：$\mathscr{L}$ 在 $C$ 上是\trAmple{}的，做法如下。先读 \S17.1.6 以理解态射 $\nu^{!?}\colon\mathit{QCoh}_C\to\mathit{QCoh}_{\widetilde{C}}$。然后，对 $C$ 上每个\trQCSheaf{} $\mathscr{F}$，描述 $C$ 上诸\trQCSheaf{}的一个正合列：
\[
0\longrightarrow\mathscr{T}\longrightarrow\nu_*\nu^{!?}\mathscr{F}\longrightarrow\mathscr{F}\longrightarrow\mathscr{Q}\longrightarrow 0,
\]
其中 $\mathscr{T}$ 与 $\mathscr{Q}$ 都支撑在一个维数为 $0$ 的集合（有限个点）上。然后证明：当 $n\gg 0$ 时，对 $i>0$ 都有 $\mathrm{H}^i(C,\mathscr{F}\otimes\mathscr{L}^{\otimes n})=0$。你将用到\trProjection{}公式（\mbox{练习~18.7.E(b)}）来证明：当 $i>0$ 且 $n\gg 0$ 时，$\mathrm{H}^i(C,(\nu_*\nu^{!?}\mathscr{F})\otimes\mathscr{L}^{\otimes n})=0$。

\textbf{18.8.2.}\ \emph{\trVeryAmple{}对\trAmple。}\quad 把「\trAmple」换成「\trVeryAmple」之后，前面的这几道练习就行不通了；这再次表明：\trAmpleness{}这个观念比\trVeryAmpleness{}表现得更规整。

\textbf{18.8.3.}\ \emph{\mbox{定理~18.8.1} 的证明。}\quad 为了得到 $(a)$--$(c)$ 蕴涵 $(e)$，利用「$\mathscr{L}^{\otimes\mathrm{N}}$ 对某个 $\mathrm{N}$ 是\trVeryAmple{}的」这一事实（\mbox{定理~16.2.2(a)}），并对 $\mathscr{F}$、$\mathscr{F}\otimes\mathscr{L}$、\dots{}与 $\mathscr{F}\otimes\mathscr{L}^{\otimes(\mathrm{N}-1)}$ 应用\trSerre{}消失（\mbox{定理~18.1.3(ii)}）。

于是现在我们假设 $(e)$，并按\mbox{定理~16.2.2} 的判别法 $(b)$ 证明 $\mathscr{L}$ 是\trAmple{}的：我们将证明，对 $X$ 上任一\trCoherentSheaf{} $\mathscr{F}$，当 $n\gg 0$ 时 $\mathscr{F}\otimes\mathscr{L}^{\otimes n}$ 都\trGloballyGenerated。

我们从一特例开始：我们将证明，当 $n\gg 0$ 时 $\mathscr{L}^{\otimes n}$ 是\trGloballyGenerated（即\trBasePointFree）。为此，只须证明：每个\trClosedPoint{} $p$ 都有一个开邻域 $U$，使得存在某个 $\mathrm{N}_p$，当 $n\ge\mathrm{N}_p$ 时，$\mathscr{L}^{\otimes n}$ 在 $U_p$ 的一切点处都\trGloballyGenerated。（理由：由\trQuasicompactness，$X$ 的每个闭子集都含有一个闭点，见\mbox{练习~5.1.E}。于是当 $p$ 跑遍 $X$ 的诸闭点时，这些 $U_p$ 覆盖 $X$。再由\trQuasicompactness，我们可以用其中有限多个 $U_p$ 覆盖 $X$。令 $\mathrm{N}$ 为相应的诸 $\mathrm{N}_p$ 中的最大者。那么当 $n\ge\mathrm{N}$ 时，$\mathscr{L}^{\otimes n}$ 在这些 $U_p$ 的每一个内都\trGloballyGenerated，从而在整个 $X$ 上亦然。）

设 $p$ 是 $X$ 的一个\trClosedPoint。对一切 $n$，$\mathscr{I}_p\otimes\mathscr{L}^{\otimes n}$ 都是\trCoh{}的（由我们的\trNoetherian{}假设）。（此处 $\mathscr{I}_p$ 是点 $p$ 的\trIdealSheaf。）由 $(e)$，存在某个 $n_0$，使得当 $n\ge n_0$ 时，$\mathrm{H}^1(X,\mathscr{I}_p\otimes\mathscr{L}^{\otimes n})=0$。由\trClosedSubscheme{}正合列
\[
0\longrightarrow\mathscr{I}_p\otimes\mathscr{L}^{\otimes n}\longrightarrow\mathscr{L}^{\otimes n}\longrightarrow\mathscr{L}^{\otimes n}|_p\longrightarrow 0,
\]
所产生的\trLongExactSequence{}可知：当 $n\ge n_0$ 时，$\mathscr{L}^{\otimes n}$ 在 $p$ 处\trGloballyGenerated。由\mbox{练习~16.1.D(b)}，存在 $p$ 的一个开邻域 $V_0$，使得 $\mathscr{L}^{\otimes n_0}$ 在 $V_0$ 的一切点处都\trGloballyGenerated。于是对一切正整数 $k$，$\mathscr{L}^{\otimes kn_0}$ 在 $V_0$ 的一切点处都\trGloballyGenerated（用\mbox{容易的练习~16.1.C}）。对每个 $i\in\{1,\dots,n_0-1\}$，存在 $p$ 的一个开邻域 $V_i$，使得 $\mathscr{L}^{\otimes(n_0+i)}$ 在 $V_i$ 的一切点处都\trGloballyGenerated（仍由\mbox{练习~16.1.D(b)}）。我们可以取每个 $V_i$ 都含于 $V_0$。由\mbox{容易的练习~16.1.C}，$\mathscr{L}^{\otimes(kn_0+n_0+i)}$ 在 $V_i$ 的每一点处都\trGloballyGenerated（因为这对 $\mathscr{L}^{\otimes kn_0}$ 与 $\mathscr{L}^{\otimes(n_0+i)}$ 都成立）。于是，在开邻域 $U_p:=\bigcap_{i=0}^{n_0-1}V_i$ 中\footnote{原书此处作『$\cap_{i=0}^{n-1}$』，当为『$\cap_{i=0}^{n_0-1}$』之误，径改。——译者注}，当 $n\ge\mathrm{N}_p:=2n_0$ 时，$\mathscr{L}^{\otimes n}$ 都\trGloballyGenerated。

到这里我们已经证明了：存在某个 $\mathrm{N}$，使得当 $n\ge\mathrm{N}$ 时，$\mathscr{L}^{\otimes n}$ 是\trGloballyGenerated。现在设 $\mathscr{F}$ 是一个\trCoherentSheaf。为了终结证明，我们将证明：当 $n\gg 0$ 时，$\mathscr{F}\otimes\mathscr{L}^{\otimes n}$ 是\trGloballyGenerated。这个论证与我们至今所做的味道相仿，所以我们把它作为练习给出。

\textbf{18.8.D.（练习）}\quad 设 $p$ 是 $X$ 的一个\trClosedPoint。

（a）证明：当 $n\gg 0$ 时，$\mathscr{F}\otimes\mathscr{L}^{\otimes n}$ 在 $p$ 处\trGloballyGenerated。

（b）证明：存在 $p$ 的一个开邻域 $U_p$，使得当 $n\gg 0$ 时，$\mathscr{F}\otimes\mathscr{L}^{\otimes n}$ 在 $U_p$ 的每一点处都\trGloballyGenerated。\emph{\trhCaveat：}虽然下面这件事是真的——由\mbox{练习~16.1.D(b)}，对每个 $n\gg 0$，都存在 $p$ 的某个开邻域 $V_n$，使得 $\mathscr{F}\otimes\mathscr{L}^{\otimes n}$ 在其上\trGloballyGenerated——但是，
\begin{equation*}\tag{18.8.3.1}\label{eq:intersection-of-neighborhoods}
\bigcap\nolimits_{n\gg 0}V_n
\end{equation*}
未必是一个开集。你可能需要利用「$\mathscr{L}^{\otimes n}$ 在 $n\ge\mathrm{N}$ 时是\trGloballyGenerated」这一事实，把 \eqref{eq:intersection-of-neighborhoods} 换成一个有限交。

\textbf{18.8.E.（练习）}\quad 证明：当 $n\gg 0$ 时，$\mathscr{F}\otimes\mathscr{L}^{\otimes n}$ 是\trGloballyGenerated，以此完成\mbox{定理~18.8.1} 的证明。\hfill$\square$

\textbf{18.8.4. \trSerre{}对\trAffineness{}的上同调刻画。}\quad \trSerre{}对\trAffineness{}的刻画类似于\mbox{定理~18.8.1}。

\textbf{18.8.5.\ \trhTheorem（\trAffineness{}的\trSerre{}上同调刻画）。}\ ——\ \emph{设 $X$ 是一个\trQuasicompact\trSeparated\trScheme。则以下诸条件等价。}

（a）\emph{概形 $X$ 是\trAffine{}的。}

（b）\emph{对 $X$ 上每个\trQCSheaf{} $\mathscr{F}$，对一切 $i>0$ 都有 $\mathrm{H}^i(X,\mathscr{F})=0$。}

（c）\emph{对 $X$ 上每个\trQCSheaf{} $\mathscr{F}$，$\mathrm{H}^1(X,\mathscr{F})=0$。}

\emph{证明。}\quad 显然 $(a)$ 蕴涵 $(b)$ 蕴涵 $(c)$（前一蕴含来自「仿射覆盖上同调消失」，\S18.1~\textbf{(vi)}），所以我们证明 $(c)$ 蕴涵 $(a)$。

令 $A=\Gamma(X,\OO_X)$，于是我们有一个典范映射 $\pi\colon X\to\Spec A$（定义于\mbox{练习~7.3.G} 及其紧随其后的那一段），并且只须证明 $\pi$ 是一个\trAffine\trMorphism。换言之，我们希望找到一些 $f_i\in A$，使得（i）诸 $f_i$ 生成 $A$（从而 $\cup\,D(f_i)=\Spec A$），而且（ii）$X_{f_i}$ 是\trAffine{}的（其中 $X_{f_i}=\{p\in X:\ f_i(p)\neq 0\}$，\mbox{定义~6.2.8}）。

我们先造出覆盖 $X$ 的一些仿射 $X_{f_i}$。对任一点 $p\in X$，取 $p$ 的一个\trAffine{}邻域 $U$，并令 $\mathscr{I}$ 为 $X\setminus U$ 的\trIdealSheaf（带诱导的既约\trSubscheme{}结构）。设 $q$ 是 $U$ 的、位于 $p$ 的闭包之中的任一\trClosedPoint（它不必是 $X$ 的闭点！）。

\textbf{18.8.F.（练习）}\quad 证明：$\mathscr{I}\to\mathscr{I}|_{\bar{q}}$ 是\trQCSheaf{}之间的一个\trSurjectiveMap，并利用假设 $(c)$ 证明它在\trGlobalSections{}上是满的。证明：存在 $f\in A$，使得 $X_f\subset U$ 是\trAffine{}的且含有 $p$。

由 $X$ 的\trQuasicompactness，我们可以用有限多个这样的 $X_f$ 覆盖 $X$，设 $f\in\{f_1,\dots,f_n\}$。为了证明 $X\to\Spec A$ 是一个\trAffine\trMorphism，我们需要 $D(f_1)\cup\cdots\cup D(f_n)=\Spec A$，即 $(f_1,\dots,f_n)=A$。

\textbf{18.8.G.（练习）}\quad 证明这一点，办法是证明 $X$ 上诸\trQCSheaf{}之间的映射
\[
\cdot(f_1,\dots,f_n)\colon\OO_X^{\oplus n}\longrightarrow\OO_X
\]
是一个满射，然后取\trGlobalSections{}并援引 $(c)$。\hfill$\square$

\textbf{18.8.H.（练习）}\quad 在与上文的\mbox{定理~18.8.5} 相同的假设下，证明 $(a)$--$(c)$ 与下述条件等价：

（d）对 $X$ 上每个\trQC{}\trIdealSheaf{} $\mathscr{I}$，$\mathrm{H}^1(X,\mathscr{I})=0$。

\textbf{18.8.6.}\quad \trSerre{}曾在\trComplexAnalyticGeometry{}中证明过一个类似的结果：\trStein{}空间同样可以用\trCoherentSheaf{}\trCohomology{}的消失来刻画。
